\section{Discussion}
\subsection{Interpretation and scope of the extension}
The principal result is a separation between simulated burden and threshold-defined extent. With identical assigned migration and growth parameters, oxygen feedback lowered viable density in the interior while leaving the lower-density boundary comparatively similar. Well-oxygenated regions retain Fisher--KPP-like growth, whereas the interior experiences inhibited births and, at sufficiently low oxygen, death transfer. The ablations show that burden-dependent supply suppression dominates the overall effect for this parameter set. The numerical comparison establishes how the selected equations behave; it does not establish that they predict glioma better than Fisher--KPP. Predictive model selection requires fitting each candidate to the same calibration observations and comparing their withheld-data predictions.

The extension is deliberately small, but its state interpretation matters. Immobile $z$ with zero clearance accumulates historical cell loss, while viable cells can diffuse into the same locations. The equations therefore allow $c+z>K$ and $z>K$ locally. A normalized necrotic volume fraction should not be inferred from $z/K$. The observed grid sensitivity reinforces this limitation. If physical packing, tissue displacement or a necrotic cavity becomes a target, a volume-fraction, mechanical or poromechanical formulation would be a different model requiring additional parameters and observations. Recent MRI-informed poromechanical forecasting and sequential biomechanical model-selection studies provide relevant alternatives, but the cited 2026 records are preprints and concern distinct preclinical datasets \citep{andaloussi2026,noman2026}. They do not validate the present closure.

\subsection{What the synthetic recovery establishes}
The two recovery experiments provide a staged inference test: homogeneous diffusion/proliferation and then separate white-/gray-matter diffusion with shared proliferation were estimated from noisy, low-resolution signals. Both include local and bootstrap intervals, independent noise trials and withheld future scans. The homogeneous generator used a finer mesh and time step than inversion. The primary BrainWeb generator instead shared the spatial grid while refining time; its separate anatomy/grid stress test exposes a limitation of that favorable design. Neither experiment removes shared biological, initialization or imaging assumptions, and neither estimates coupled oxygen--necrosis mechanisms from structural MRI.

The homogeneous known rates lay within both interval constructions. In the BrainWeb primary realization, the true proliferation rate narrowly missed both 95\% intervals, while both diffusion rates were contained. Retaining that outcome is essential: a nominal confidence level does not imply containment in every realized dataset. Independent local-interval coverage of \BWWhiteCoverage/\BWTrials{}, \BWGrayCoverage/\BWTrials{} and \BWRhoCoverage/\BWTrials{} had broad binomial uncertainty. These observations do not establish bootstrap coverage, global identifiability or calibration across patients. The weaker relative precision of gray-matter diffusion and its dependence on white-matter diffusion argue for propagating the joint rate uncertainty. Separate tissue rates were estimable here without fixing their ratio, but that conclusion depends on this known seed, informative tissue exposure and observation schedule. Relaxing initialization and calibration may expose stronger confounding \citep{hiremath2025}.

The anatomy test identifies a distinct source of uncertainty. Noise-free 2 mm signals fitted with the 3 mm model shifted gray-matter diffusion by +\BWGrayStressShift\%, despite day-90 observation discrepancy remaining below the noise standard deviation. Many measurements can make sampling intervals narrow while systematic representation error persists. This shift combines tissue-boundary resampling, mesh, seed sampling and discrete imaging; it is not a measured biological change, a pure spatial-order result or an estimated segmentation-error distribution. The primary bootstrap and local covariance keep these quantities fixed and cannot quantify this effect. Inference should therefore compare multiple anatomically consistent resolutions and propagate plausible label/registration variations before interpreting tissue-rate differences biologically.

Forecast assessment also needs spatial qualification. In BrainWeb, day-90 prediction intervals contained 94.2\% of full-FOV measurements but 89.2\% near the seed. Background-inclusive summaries alone can conceal local mismatch. Mean-signal confidence bands and new-measurement prediction bands answer different questions; neither provides simultaneous image coverage here. These frequentist conditional intervals are not Bayesian posteriors or estimates of clinical variability. Next studies should vary seed location, tissue-rate contrast and scan timing, and introduce unknown age, seed shape, capacity, calibration, correlated noise and structural mismatch. A broader bootstrap-coverage audit remains separate work.

\subsection{What measurements could identify parameters?}
A practical first stage would estimate a limited growth--migration subset from serial, registered MRI, while jointly assessing initialization and whether separate tissue migration rates are supported. A constrained tissue ratio or an external prior may be needed when the measurements do not separate those rates; the conditional BrainWeb example does not establish that this restriction can be removed in patients. In a homogeneous, isotropic Fisher--KPP medium, the asymptotic minimal front speed is $v=2\sqrt{D\rho}$ and the front length scale is proportional to $L=\sqrt{D/\rho}$. These relations explain why growth velocity and spatial spread offer complementary information. They are not identities for a finite, heterogeneous brain with an unknown initial state. One snapshot with unknown tumor age generally cannot separate all dimensional rates. Recent single-scan inference methods introduce model and scaling assumptions \citep{zhang2025}; a 2026 simulated-biopsy preprint similarly distinguishes information about infiltration length from information about growth speed \citep{hofmann2026}.

\begin{table}[p]
\centering\small
\caption{Proposed measurement and prior strategy. These are candidate information sources, not measurements acquired or estimates obtained in this study. Each modality requires its own observation operator and uncertainty model.}\label{tab:measurements}
\begin{tabularx}{\textwidth}{p{0.17\textwidth}YY}
\toprule Quantity & Potential information & Restriction or prior needed\\\midrule
$c(0)$, $D$, $\rho$ & Serial T1/T2/FLAIR, diffusion-weighted and quantitative MRI; registered anatomy & Joint initialization uncertainty; estimate a small parameter subset; constrain tissue ratios. MRI water diffusion is not cell migration.\\
$K$ and cell scale & Spatially localized tumor-cell counts with tissue-volume or stereological calibration & Distinguish tumor from non-tumor cells and correct sampling/deformation; MRI alone does not fix the absolute scale.\\
$\beta$, $\alpha$ & Oxygen-sensitive imaging combined with perfusion/vascular measurements and externally calibrated uptake information & Supply and consumption are confounded; blood-flow surrogates are not the restoration coefficient itself.\\
$\ell_o$ & External transport measurements or multiscale vascular/oxygen models & Constrain externally when subvoxel; the current 1--2 mm meshes do not resolve 0.2 mm.\\
$\theta_p$, $\theta_d$, $\delta$ & Controlled oxygen-dependent growth and survival experiments; spatially matched hypoxia/viability observations & Relative-oxygen scale must be calibrated. Ki-67 is a cycling fraction, not directly a growth rate; hypoxia markers are not absolute oxygen readings.\\
$\lambda_z$ and $z$ mapping & Longitudinal necrotic-region imaging plus localized pathology & Separate formation from clearance and cell equivalents from tissue volume; a single scan cannot determine both histories.\\\bottomrule
\end{tabularx}
\end{table}

MRI-derived cell-density maps can supply more information than binary segmentations only when their calibration is defensible. ADC measures water mobility, not the effective migration coefficient $D$. Its association with cellularity is influenced by extracellular space, edema, tissue composition and necrosis \citep{eidel2016}. Nor should a FLAIR boundary automatically be treated as a fixed tumor-density contour: an MRI--histology study challenged that assumption in a single autopsy case \citep{martens2021}. In a 2025 study of image-guided biopsies, associations between MRI-derived cell-density maps and histology differed between untreated and recurrent tumors \citep{nocera2025}. The treatment context and observation mapping must therefore be part of the inference model. Biopsy coordinates, brain shift, fixation deformation, partial volume and the conversion from two-dimensional counts to cells/mm$^3$ require explicit treatment.

Oxygen introduces further observation ambiguity. The state $o$ is availability relative to a normal baseline; it is not directly oxygen extraction fraction, hemoglobin saturation or oxygen tension. Oxygen-enhanced MRI, BOLD-based measures and hypoxia PET could constrain an observation model, but their signals depend on delivery, vascular and tissue factors. A histological comparison of MRI-derived oxygen measures illustrates that an imaging surrogate cannot simply be equated with a hypoxia marker \citep{arzanforoosh2024}. The 2026 glioma habitat repeatability study is encouraging for measurement consistency, but repeatability does not establish an absolute oxygen calibration or biological specificity \citep{das2026}.

Priors are necessary where these measurements leave parameters weakly identified. Positive rates and lengths can use log-normal priors; ordered oxygen thresholds can use a transformed prior satisfying $0<\theta_d<\theta_p<1$. Hyperparameters should come from independently evaluated experiments or a defensible population model, with uncertainty reflecting transfer between systems. The working values in Table~\ref{tab:parameters} are not empirical prior means. Prior predictive checks should reveal implausible growth, oxygen depletion and accumulation above the intended interpretation of $K$. Sensitivity analysis, profile likelihoods and posterior correlations can then identify parameter combinations that the observations fail to distinguish. The 2025 identifiability/model-selection study provides a relevant framework, although its virtual-cohort chemoradiation models differ from this system \citep{hiremath2025}. A narrow posterior produced mainly by a strong prior should not be described as precise measurement-based estimation.

\subsection{Multimodal MRI and patient-specific growth forecasts}
Serial T1, gadolinium-enhanced T1 (T1-Gd), T2/FLAIR, diffusion MRI (dMRI) and functional MRI (fMRI) could constrain complementary parts of a patient-specific model (Table~\ref{tab:multimodal}). Structural images would define anatomy, abnormal tissue and the postoperative cavity; diffusion-derived microstructure could inform migration geometry and a calibrated cellularity observation model. Enhancement reflects blood--brain barrier permeability as well as tumor biology, while FLAIR abnormality can combine infiltration, edema and treatment-related changes \citep{wen2023,martens2021}. Thus neither boundary should be imposed as a universal cell-density contour. The synthetic normalized-cellularity operator used here is not an established mapping from any of these clinical sequences.

\begin{table}[tbp]
\centering\small
\caption{Proposed roles of multimodal MRI in a growth--brain digital twin. The PAT16 pilot uses existing T1, dMRI and resting-state fMRI derivatives for anatomical/directional integration and static connectivity; no patient growth rates are fitted. T1-Gd and T2/FLAIR integration remains proposed. Imaging correlates require calibration and do not directly identify every model state.}\label{tab:multimodal}
\begin{tabularx}{\textwidth}{>{\raggedright\arraybackslash}p{0.13\textwidth}YY}
\toprule Modality & Information for a proposed model & Interpretation and uncertainty\\\midrule
T1 & Anatomy, tissue compartments, displacement and resection cavity & Lesion-aware segmentation and registration; non-enhancing margins may be poorly defined.\\
T1-Gd & Enhancing compartment and permeability-related change & Enhancement is not a direct viable-cell count; surgery and treatment alter permeability.\\
T2/FLAIR & Extent of abnormal tissue and non-enhancing disease & Infiltration, edema, gliosis and treatment effects can overlap.\\
dMRI & Fiber orientations, tract architecture and diffusion-derived microstructure & Water diffusion is not cell migration; ADC-to-cellularity calibration is context dependent.\\
Resting-state fMRI & Functional connectivity and hemodynamic constraints for neural-network models & BOLD can be altered by neurovascular uncoupling; a short time series does not measure longitudinal tumor growth.\\\bottomrule
\end{tabularx}
\end{table}

The PAT16 pilot implements positive-definite, trace-matched white-matter migration tensors informed by dMRI orientations (Equation~\eqref{eq:pat16_tensor}). Parallel and transverse migration scales, and their relation to gray-matter migration, would remain parameters or externally constrained quantities; the measured water-diffusion tensor must not be substituted numerically for tumor-cell migration. Crossing fibers, edema, tract displacement and loss introduce uncertainty in the directional information itself. The scalar implementation remains a matched reference in the pilot; predictive patient comparison would require separate calibration of both variants. The anatomy sensitivity observed here argues for evaluating plausible segmentations, registration transforms and directional reconstructions jointly with rate uncertainty before attributing a fitted difference to tumor biology.

More MRI contrasts do not automatically identify the oxygen--necrosis subsystem. Perfusion measurements may help constrain delivery, while dedicated oxygen-sensitive MRI or hypoxia PET may constrain a calibrated relation to $o$; ordinary resting-state BOLD is not a tissue oxygen-concentration measurement. The 2026 repeatability work used a dedicated multiparametric protocol in a murine glioma model, rather than validating oxygen estimates from routine human fMRI \citep{das2026}. Likewise, a postoperative cavity or a radiological necrotic region cannot be equated with the accumulated cell equivalents $z$. Observation models may need additional edema, vascular or treatment-injury quantities that are absent from the present equations.

\subsection{Connecting growth, brain networks and functional outcomes}
Growth forecasts could provide the evolving lesion state for a brain digital twin. Previous studies used individual structural connectomes and The Virtual Brain to model functional connectivity in tumor patients and after resection \citep{aerts2018,aerts2020}. The related Finse 2026 repository demonstrates a T1/dMRI/resting-state-fMRI workflow linking anatomy, biomechanics, connectivity, neural dynamics and cognitive outcomes in one example, sub-PAT16 \citep{lundervold2026twin}. The present PAT16 pilot reads its processed MRI derivatives and recomputes static functional connectivity, adding a conditional migration comparison. It does not couple growth to mechanics, evolving structural connectivity or neural dynamics; those remain proposed extensions.

The underlying Brain Tumor Connectomics (BTC) data comprise 11 glioma patients, 14 meningioma patients and 11 controls before surgery, with postoperative data available for 7, 12 and 10 participants, respectively, around six months later \citep{aerts2022}. Gliomas and meningiomas should not be pooled under the same infiltrative growth law. The linked prototype supplies T1, dMRI and resting-state fMRI; T1-Gd, T2 and FLAIR would be additions. A preoperative/six-month-postoperative pair combines resection, recovery and any intervening treatment with disease evolution. It cannot be interpreted as an untreated growth interval or independently determine all migration, proliferation and response parameters. Lesion-aware registration must distinguish removed tissue from displacement; a cavity-crossing deformation is not a measured growth trajectory. Serial clinical tumor MRI, recorded treatment timing, molecular diagnosis and an explicit residual-disease state would be needed for a defensible growth-inference study.

A staged integration would first propagate ensembles of plausible tumor states through mechanics and tract-involvement models, then map these changes to structural connectivity and neural dynamics. Both direct infiltration and tissue displacement would be candidate causes of altered connectivity, with distinct parameters and observations. Resting-state fMRI would constrain the functional/hemodynamic component, while later functional measurements and cognitive scores would test its forecasts. Tumor growth over days or months and BOLD fluctuations over minutes are different time scales. Neurovascular uncoupling can produce apparent loss of functional connectivity without corresponding loss of neuronal function \citep{sun2020}; weak BOLD must therefore not be treated as proof that tissue is functionally expendable. Fitting a connectivity matrix and evaluating agreement with that same matrix is calibration, not independent validation of functional prognosis.

Bidirectional neuron--glioma coupling is a further hypothesis. Human glioblastoma work combining electrophysiology, localized tissue sampling and experimental models provides evidence for tumor--neural circuit interactions \citep{krishna2023}. It does not establish a transferable relation between resting-state BOLD amplitude or connectivity and proliferation rate, particularly across glioma subtypes. An activity-dependent growth term would require its own observation mapping, identifiable coupling parameters and independent longitudinal or perturbation evidence. Initial integration should assess functional consequences of growth before attempting to infer this feedback from sparse imaging.

For treatment evaluation, an integrated model could compare resection or radiation scenarios in terms of both residual/infiltrative burden and predicted tract, network and cognitive consequences. A next PAT16 study could extend the static overlays to conditional network effects using the observed resection and anatomical uncertainty, without claiming recovery of untreated growth or the unobserved outcome of alternative surgery. Later studies should compare growth-only, network-only and combined predictors on patient-level withheld outcomes; retrospective imaging or connectivity agreement does not establish benefit from changing treatment. Functional preservation would join tumor control, toxicity and deliverability as a separately validated objective.

\subsection{Interpretation of the PAT16 directional pilot}
The paired comparison isolates a transport-geometry choice under shared initialization, local mean diffusivity and proliferation, rather than estimating different patient growth rates. The broad marginal sensitivity bands overlap, while pairing reveals a negative day-90 burden difference for the recorded draws. This direction is conditional on the lesion, coefficients, interfaces and initialization. Trace matching fixes $\operatorname{tr}(\mathbf D)/3$ but does not impose equal spatial spreading in a bounded heterogeneous domain. Diffusion conserves mass; its redistribution changes local saturation and therefore integrated logistic births. The anisotropic calculation contains no cell-killing term, and the smaller burden is not evidence of a therapeutic effect or universally slower glioma growth. Regional redistribution can also occur despite a lower global burden.

The pilot moves beyond a healthy anatomical phantom by using a measured lesion, tissue maps and tract axes, but the state-to-image link remains assumed: a lesion mask is neither a measured cellularity field nor a calibrated MRI detection contour. A 3 mm staircase interface and the treatment of 603 lesion/CSF-classified voxels influence the transport graph. The small time-step difference and analytical tensor tests do not establish spatial convergence on this anatomy. DTI cannot resolve crossing fibers; the existing multi-shell tensor fit, FA eligibility rule, segmentation and registration require further assessment. A global 3-degree direction perturbation is a limited sensitivity probe, not a registration-error posterior. The 32-draw, independently specified ensemble omits correlated and structural uncertainty and cannot support calibrated patient risk estimates.

The static fMRI overlays locate simulated tissue involvement relative to measured correlations. They demonstrate coordinate integration and offer hypotheses for later network experiments, while regional mean density remains an anatomical overlap measure. FC edges are not axonal pathways, and projecting ROI centers through a slice slab does not locate a tract on that plane. Motion, atlas mismatch and neurovascular uncoupling affect the measured FC. Independent longitudinal connectivity and cognitive measurements, with explicit neural and hemodynamic models, are required before interpreting tumor overlap as functional impairment.

\subsection{Uncertainty quantification and informative measurements}
The positive rate correlation and forecast ensembles in the synthetic experiment illustrate why parameter estimation should quantify joint uncertainty rather than supply a single best fit. Beyond the conditional frequentist analysis performed here, a Bayesian analysis could combine independently justified priors with an observation likelihood accounting for imaging noise and spatial correlation. Hierarchical population priors could support patients with sparse imaging, provided the population model is learned within training data. Uncertain initialization, registration and segmentation should be included as nuisance quantities or evaluated through repeated analyses. Strong posterior correlations and sensitivity to alternative priors would indicate which rates or combinations remain poorly constrained.

Forecast ensembles should propagate these uncertainties into viable-burden intervals, infiltration maps and oxygen/necrosis predictions. Measurement uncertainty, numerical error and model inadequacy require separate assessment: a narrow conditional posterior does not account automatically for a wrong oxygen closure or unresolved spatial gradients. Competing Fisher--KPP and coupled formulations should be evaluated separately on withheld observations; any model-averaging weights should be learned without those evaluation data. Posterior predictive checks, interval coverage and width, and sensitivity to the observation operator should accompany point-error metrics.

Pash et al. developed a scalable Bayesian reaction--diffusion framework using longitudinal imaging and investigated imaging frequency and experimental design \citep{pash2026}. This provides a methodological reference for the proposed inference workflow, not validation of the present oxygen--necrosis extension. Once uncertainty is calibrated, candidate scan times or additional modalities could be ranked by their expected reduction in forecast uncertainty or their expected value for distinguishing treatment options. Additional measurements are useful when they resolve a consequential ambiguity, rather than merely increasing the number of fitted parameters.

\subsection{A validation pathway with independent targets}
Validation must be tied to a specified observable, population and prediction horizon. Five complementary stages are appropriate here.
\begin{enumerate}
\item \textbf{Verify the solver.} Check conservation, positivity, analytical limits, solver tolerances and time/space refinement. These checks address mathematical implementation. The present unresolved oxygen length and necrotic peak require further spatial assessment before local quantitative claims.
\item \textbf{Extend synthetic recovery.} The homogeneous and BrainWeb experiments supply controlled two- and three-rate recovery and withheld-scan checks. Broaden them across seed locations, rates, anatomy representations and scan schedules, and add unknown initialization, correlated noise, registration error and observation-model mismatch. Use a different mesh, solver or observation formulation to reduce the inverse crime of fitting data generated by the identical numerical pipeline. Synthetic success does not establish biological adequacy; earlier MRI-constrained reaction--diffusion work demonstrated the importance of separating synthetic recovery from in vivo forecasting \citep{hormuth2015}.
\item \textbf{Forecast withheld observations.} Fit early observations, freeze the calibration and measurement mapping, and predict later scans over a declared horizon. Compare separately calibrated M0, M1 and M2 with a simple growth baseline. Treated tumors require explicit treatment terms or a justified untreated interval. Updating with each new scan is sequential assimilation, a different evaluation from a frozen forecast. The 2025 assimilation study illustrates the distinction in high-grade glioma during chemoradiation \citep{miniere2025}.
\item \textbf{Check latent fields independently.} Reserve spatially localized cell counts, viable/necrotic histology and oxygen-related observations for evaluation rather than using all samples for calibration. Report localization and measurement uncertainty. A successful volume forecast alone does not validate $o$, $z$ or the death mechanism. Controlled preclinical oxygen or perfusion perturbations would provide stronger tests of the proposed feedback than morphology alone.
\item \textbf{Evaluate transportability.} Apply a fixed analysis and prior strategy to an external cohort, with any permitted patient-specific calibration specified in advance. Group splits by patient and, where relevant, site and treatment status. Prospective clinical utility is a subsequent question, beyond retrospective image agreement.
\end{enumerate}

Evaluation should include volume errors, surface distances and overlap, together with quantitative field errors where calibrated measurements exist. Report predictive interval coverage and width, not only point accuracy. Repeated biopsies within a patient and spatially correlated voxels are not independent replicates; patient-level or hierarchical analyses should reflect this dependence. The 2025 biopsy-validation study is a useful example of independent tissue evaluation and context-dependent agreement \citep{nocera2025}. An accurate forecast can coexist with non-unique physiological parameters, while an identifiable model can still have the wrong biology. Both prediction and parameter interpretation must be assessed.

Clinical validation should subsequently use a prespecified population, prediction horizon and analysis, with external testing before prospective, blinded forecasting. Imaging endpoints should account for treatment effects and pseudoprogression using a defined response-assessment framework such as RANO 2.0 \citep{wen2023}. Agreement with a later MRI or recurrence location is evidence about prediction under the observed treatment; it does not establish the outcome under an alternative treatment. Demonstrating clinical utility would require prospective comparative evaluation of decisions and, ultimately, controlled trials assessing survival, neurological function, quality of life and toxicity.

\subsection{Treatment evaluation and candidate strategies}
The present simulations suggest why treatment evaluation should consider viable burden and spatial infiltration together: similar threshold-defined volumes can conceal different modeled interiors. Their difference is not a treatment effect, because no therapy was simulated. Extending the framework would require an explicit resection geometry and residual disease state, spatially delivered radiation and drug exposure/response terms. Surgery may also change anatomy and perfusion, requiring updated transport and observation models. Natural growth and treatment-response parameters should be estimated jointly from appropriately timed observations, or constrained independently when they cannot be separated.

An oxygen-dependent radiation-survival model could test the consequences of hypoxia and reoxygenation between fractions, but the relative availability $o$ first needs a calibrated connection to the relevant radiobiology. Neither $c$ nor $z$ currently represents a measured clonogenic population or radiation-induced normal-tissue injury. Computing a tumor-control probability would require additional assumptions about clonogenic fraction and survival; evaluating toxicity would require a separate normal-tissue model. Treatment-related necrosis and imaging changes must be distinguished from spontaneous hypoxic loss. These additions are proposed mechanisms requiring independent experiments, not interpretations already supported by the three simulated fields.

Recent studies illustrate complementary evaluation targets. Balcerak et al. evaluated imaging- and physics-informed infiltration estimates against subsequent recurrence in a dataset of 152 glioblastoma patients \citep{balcerak2025}. H\"ager et al. examined model-informed target delineation and voxel-based dose planning in three cases, demonstrating planning feasibility while identifying difficulties in reproducing complex dose distributions \citep{hager2026}. These retrospective prediction and planning studies do not establish a survival benefit from delivering the proposed alternatives. Earlier risk-aware radiotherapy optimization used a cohort of 100 virtual patients \citep{chaudhuri2023}, illustrating how uncertainty can inform simulated treatment comparisons without constituting a clinical trial.

A next experimental stage could use paired virtual cohorts to compare a specified standard-treatment comparator with adaptive fractionation, spatial dose painting, hypoxia-modifying combinations and alternative drug/radiation sequences. Each candidate would require an explicit mechanism and independently constrained response parameters. Comparisons should use matched dose/resource constraints, assess normal-brain exposure and delivery feasibility, and test whether a favorable ranking persists across parameter and model uncertainty. Adaptive policies must use only observations available at the decision time. Outcomes could include viable burden, spatial recurrence and time to a prespecified progression criterion; a survival prediction would require a separately validated link to that endpoint. Controlled preclinical oxygen/perfusion perturbations could test the proposed feedback before patient translation. Virtual cohorts should prioritize experimentally testable hypotheses and informative measurements; they cannot identify an effective new therapy solely from a favorable simulated result.

\subsection{Limitations and staged roadmap}
The anatomical comparison uses one healthy phantom, one seed and one primary parameter set. The ablations and refinement runs are not a population analysis. The recovery experiments vary noise around two fixed states, including one BrainWeb seed; they do not sample between-patient heterogeneity or calibrate the dense-seed oxygen--necrosis comparison. The 3 mm primary BrainWeb fit shares tissue geometry with its generator, and one 2 mm sensitivity example is not a distributional audit of anatomy uncertainty. Scalar migration in the BrainWeb comparison omits fiber orientation; the separate single-patient pilot adds DTI-informed transport but no spatial convergence or patient growth validation. No displacement, edema, vascular evolution, immune interactions, phenotype switching or therapy is modeled. Quasi-steady oxygen identifies effective coefficient groups rather than absolute oxygen equilibration rates. Minimum oxygen and necrotic accumulation are under-resolved, and zero clearance accentuates accumulation. Finer simulation grids do not create additional measurement information.

The following sequence makes the roadmap concrete, building on the numerical comparisons and conditional synthetic recovery reported here.
\begin{enumerate}
\item \textbf{Establish an estimable MRI model.} Resolve the coupled-model numerical limitations and calibrate sequence-specific T1/T1-Gd/T2/FLAIR and diffusion observation operators. Extend two- and three-rate recovery across uncertain initialization and anatomy and extend the PAT16 directional comparison to separately fitted isotropic and anisotropic models. Add oxygen/death parameters only when independent measurements or justified priors constrain them.
\item \textbf{Broaden uncertainty assessment.} Propagate initialization, registration, tissue-label, tract-direction and imaging-model uncertainty, explicitly testing the anatomy sensitivity found here. Audit coverage, bootstrap calibration and prior sensitivity while separating numerical and model errors.
\item \textbf{Demonstrate external forecasting.} Fit early scans and predict withheld later observations with separately calibrated variants and simple baselines. Distinguish untreated growth from postoperative/treatment response, then test an external cohort and independent spatial biological targets.
\item \textbf{Connect growth to functional outcomes.} Build on the PAT16 anatomical integration and static overlays, then test mechanics/tract/connectome couplings and neural dynamics against withheld fMRI and cognitive outcomes across patients. Preserve a separate hemodynamic observation model and test whether integration improves on growth-only and network-only baselines.
\item \textbf{Evaluate treatment hypotheses.} Implement and verify treatment and normal-tissue effects, calibrate their response parameters, and compare feasible strategies with matched treatment constraints in virtual cohorts and controlled experiments. Include functional preservation and informative imaging schedules.
\item \textbf{Test clinical utility.} Evaluate prospective, blinded forecasts before controlled trials of model-guided care, with prespecified patient benefit and harm endpoints.
\end{enumerate}
Add complexity only to address predictive discrepancies supported by independent observations. Each stage requires its own evidence.

\subsection{Conclusion}
A minimal oxygen--necrosis extension can be implemented alongside Fisher--KPP within the same anatomical and interactive workflow. Under the assigned conditions, the extension reduced viable burden much more than low-threshold extent, largely through growth inhibition. Its localized necrotic predictions were sensitive to grid spacing and cannot be interpreted as validated tissue fractions. Staged synthetic experiments estimated homogeneous and tissue-specific baseline rates with conditional uncertainty and withheld forecasts. A proliferation interval miss was retained, and anatomy/grid changes produced a material gray-matter rate shift that noise-based intervals do not capture. The PAT16 pilot additionally demonstrates trace-matched directional growth and static MRI/network integration under assumed parameter distributions. Together, the code, recorded experiments and manuscript provide a tested starting point for measurement-based inference and staged validation; patient calibration and independent biological evaluation remain necessary.
